% Independent bounded algebra check; additive fragment.
\subsection{Reflection and the original resolvent: a direct verification}

This verifies NCE10--13 on the exact objects of EC
\cite{prog:EC}, and also records the orientation of NCE3--9.
All inner products are conjugate linear in the first entry.
Put \(c=a/2\).  The coordinate change used solely to express physical
reflection is \(S=c+iy\), with no change of the source mass.

The roots of the monic polynomial
\(\widehat\chi(y)=i^{-q}\chi(c+iy)\), retaining their complete
multiplicities, are
\[
 y_{rt}=(2t-a)\gamma-i(2r-a)\delta .
\]
Conjugation permutes these roots by \(r\mapsto a-r\), while
negation permutes them by \((r,t)\mapsto(a-r,a-t)\).
Since the polynomial is monic, these two exact permutations give
\begin{equation}
 \overline{\widehat\chi(\bar y)}=\widehat\chi(y),\qquad
 \widehat\chi(-y)=(-1)^q\widehat\chi(y),\qquad
 \chi(a-S)=(-1)^q\chi(S).
 \tag{ECR1}
\end{equation}
Thus the physical polynomial is real, with parity \((-1)^q\).
The last identity, including when \(q\) is odd, implies that
\(\mathscr R f(S):=f(a-S)\) preserves
\(\chi\mathcal P_{N-q}\).

Evenness of the original measure makes \(\mathscr R\) a unitary
involution.  It preserves \(\mathcal P_N\), and hence preserves
\(\mathcal V_N=\mathcal P_N\ominus\chi\mathcal P_{N-q}\) and commutes
with its orthogonal projector \(\Pi_N\).
The latter statement follows directly from the unique orthogonal
decomposition into \(\mathcal V_N\) and its orthogonal complement,
both of which are invariant under the unitary involution.
The attained section \(T_N\) consequently intertwines reflection on
classes with reflection on representatives.

For the monic original source polynomial \(p_n\), the polynomial
\((-1)^np_n(a-S)\) is monic of degree \(n\).
Its inner product with any polynomial of smaller degree is zero by
unitarity and the defining orthogonality of \(p_n\).
Uniqueness of the monic orthogonal polynomial follows because the
difference of two candidates has lower degree and is orthogonal to
itself.  Therefore
\begin{equation}
 p_n(a-S)=(-1)^np_n(S),\qquad
 \mathscr R\widehat b_n=(-1)^n\widehat b_n,\qquad
 \langle\widehat b_N,\widehat b_{N+1}\rangle=0 .
 \tag{ECR2}
\end{equation}
Here \(b_n=[p_n]\), \(\widehat b_n=T_Nb_n\), and the same fixed
attained section is used for both \(n=N,N+1\).
The last equality follows by applying the unitary involution to both
arguments: it negates their inner product.

The operator \(J=\Pi_NM_y|_{\mathcal V_N}\) satisfies
\(\langle f,Jg\rangle=\langle Jf,g\rangle\), because \(y\) is real
in the original measure.  Also
\(\mathscr RM_y=-M_y\mathscr R\), and commutation of \(\Pi_N\)
with reflection therefore proves
\(\mathscr RJ=-J\mathscr R\).
These identities are valid for either parity of \(q\).

For \(f=T_Nv\), the polynomial
\(Sf-(\ell_Nv)p_{N+1}\) has degree at most \(N\), because \(p_{N+1}\)
is monic and \(\ell_Nv\) is the literal leading coefficient of \(f\).
Its class is \(Av-(\ell_Nv)b_{N+1}\).  Projecting it to
\(\mathcal V_N\) thus gives the attained representative of that
class.  Since \(p_{N+1}\) is orthogonal to every polynomial of degree
at most \(N\), the other evaluation of this same projection is
\(\Pi_NSf=cf+iJf\).  Finally,
\(\ell_Nv=\langle\widehat b_N,f\rangle/\omega_N\), where
\(\omega_N=\|p_N\|^2\).  Hence
\begin{equation}
 T_NAT_N^{-1}f
 =cf+iJf+
   \widehat b_{N+1}\frac{\langle\widehat b_N,f\rangle}{\omega_N}.
 \tag{ECR3}
\end{equation}
This proves the positive sign in NCE10 on its actual domain.

Write \(\lambda=c+d+it\), with \(d=a\delta>0\), \(t=a\gamma>0\),
and let \(v=v_\lambda\) be the original retained eigenclass.
Put \(f=T_Nv\), \(h=\widehat b_{N+1}\),
\(F_{13}=\langle\widehat b_N,f\rangle\),
\(F_{23}=\langle h,f\rangle\), and \(E=\|f\|^2\).
The leading-row and boundary-column identifications are, directly,
\[
 \ell_N=\omega_N^{-1}b_N^*G_N,\qquad \xi_N=-G_Nb_{N+1}.
\]
Indeed, \(p_N-T_Nb_N\) is an old relation, and
\(\langle p_N,T_Nx\rangle=\omega_N\ell_Nx\).
Likewise \(p_{N+1}-\chi S^{N+1-q}\) has degree at most \(N\),
represents \(b_{N+1}\), and its pairing with \(T_Nx\) equals
\(-\langle T_Nx,\chi S^{N+1-q}\rangle\).
These prove the two identifications without exchanging the row
orientation.  Pairing \((S-\lambda)f\) with \(f\) now gives
\[
 -\frac{\overline{F_{23}}F_{13}}{\omega_N}
   =E\{-d+i(\mu-t)\},\qquad
 \mu=\frac{\langle f,Jf\rangle}{E}.
\]
Conjugation and negation prove
\begin{equation}
 \mathfrak c:=\frac{\overline{F_{13}}F_{23}}{\omega_NE}
     =d+i(\mu-t).
 \tag{ECR4}
\end{equation}
In particular both Gram entries are nonzero since \(d>0\).

The Hermitian property of \(J\) proves
\(\operatorname{Re}\langle x,(d+it-iJ)x\rangle=d\|x\|^2\).
Thus this finite square operator is invertible.  Let
\(R=(d+it-iJ)^{-1}\).  Applying (ECR3) to the eigenclass gives
\(f=(F_{13}/\omega_N)Rh\).  Substitution into \(F_{23}\) and \(E\)
then yields, with its stated unconjugated numerator,
\begin{equation}
 \mathfrak c
  =\frac{\langle h,Rh\rangle}{\|Rh\|^2},\qquad
 \mu=\frac{\langle Rh,JRh\rangle}{\|Rh\|^2}.
 \tag{ECR5}
\end{equation}
These are exact identities in the original source Hilbert space.

Diagonalize the finite Hermitian operator \(J\); existence of its
orthogonal eigenbasis follows by the elementary Rayleigh-quotient
argument supplied in NAB13.
Reflection sends its \(x\)-eigenspace unitarily onto its
\((-x)\)-eigenspace, and (ECR2) preserves \(h\) up to a sign.
The squared lengths of the two spectral projections of \(h\) are
therefore equal.  For a pair \(x,-x\), the mean under the positive
weights introduced by \(\|Rh\|^2\) is exactly
\begin{equation}
 \frac{x\{[d^2+(t-x)^2]^{-1}-[d^2+(t+x)^2]^{-1}\}}
 {[d^2+(t-x)^2]^{-1}+[d^2+(t+x)^2]^{-1}}
   =\frac{2tx^2}{d^2+t^2+x^2}.
 \tag{ECR6}
\end{equation}
It lies in \([0,2t)\), and is strictly positive for \(x\ne0\).
The spectral measure cannot be supported entirely at zero: that
would give \(Rh=(d+it)^{-1}h\), hence \(f\) proportional to \(h\),
which by (ECR2) implies \(F_{13}=0\), contradicting (ECR4).
The total mean in (ECR5) is a positive weighted average of the
paired means, with some nonzero spectral mass away from zero.
Consequently \(0<\mu<2t\),
\(|\operatorname{Im}\mathfrak c|<t\), and
\(|\mathfrak c|<\sqrt{d^2+t^2}\), proving NCE12--13 with their
strict signs.

The accompanying exact test uses \(a=2,\delta=1/4,\gamma=3\),
\(q=9,N=9\), and a Gaussian measure.  It retains the original
\(S=1+iy\) coefficients and a nonzero old-relation space, and checks
the odd-\(q\) reflection, both boundary parities, (ECR3), (ECR4),
and (ECR5) by exact rational arithmetic.
This is an algebraic diagnostic only, not a packet, period, or
arithmetic-source evaluation.
